\scanpage{36}
17) \uline{Relations entre classes fond locales et classes fondamentales globales}

\diagram{\begin{tikzcd}[column sep=large,row sep=large]
Y \arrow[r,hook,no head,"i","(d)"'] \arrow[d,"g\ (r\deleted{\unsure{+d}})"',"\deleted{\unsure{f}}"] & X \arrow[dl,no head,"\unsure{f|g}\ (r+d)"] \\
S &
\end{tikzcd}}

\deleted{\[
\begin{aligned}
T_{Y/X}\otimes T_{Y/S}&\simeq T_{X/S}|Y\\
T_{Y/S}&\simeq\underline H^{2d}_Y(T_{X/S})\\
R^{2r}_!f(\deleted{\unsure{g}}T_{Y/S})&\simeq\\
T_{Y/S}&\simeq\unsure{+}
\end{aligned}
\]}

\[
\left\{\begin{aligned}
\underline H^{\deleted{\unsure{2d}}\inserted{i}}_Y(T_{X/S})&=0&&\text{si }i\ne2d\\
\underline H^{2d}_Y(\unsure{T|G}_{X/S})&\simeq T_{Y/S}
\end{aligned}\right.
\]

\diagram{\begin{tikzcd}[column sep=small,row sep=large]
\displaystyle\Rightarrow\left\{\begin{array}{l}R^{\deleted{2d}\inserted{i}}_!g_{*Y}(T_{X/S})=0\quad\text{si }i<2d\\[3em]R^{2r}_!g_{*Y}(T_{X/S})\end{array}\right. \arrow[r,phantom,"\simeq" description] & R^{2r\deleted{\unsure{+2d}}}_{\unsure{!}}g_{*Y}(\underline H^{2d}_Y(T_{X/S})) \arrow[r,phantom,"\simeq" description] \arrow[d] & R^{2r+2d}_{\substack{\text{supports}\\\text{\unsure{pr. sur}}/S\\\text{et }\subset Y}}g_*(T_{X/S} \\
\deleted{\unsure{H^0}} & R^{2r+2d}_!g_{*X}(T_{X/S})) \arrow[dl] & \\
(\mathbf Z/n\mathbf Z)_S & &
\end{tikzcd}}

Compatibilité avec les hom traces ....

S'exprime plus aisément quand on dispose de la catégorie dérivée, et s'énonce sous forme plus complète.

18) \uline{Catégories dérivées} et application aux $R_!f_*$.
\figure{\begin{tikzpicture}[scale=.45]\draw (0,1.1)--(.23,.65)..controls(.58,.4)and(.4,.25)..(.08,.48)..controls(-.12,.65)and(.4,.78)..(.36,.89)--(.05,.91);\draw (.21,.65)..controls(.7,.15)and(.45,-.3)..(.24,-.55)--(-.05,-.8);\draw(.24,-.55)--(.45,-.65);\end{tikzpicture}}

19) \uline{Th. de dualité globale}.
\mspage{35}
